\subsection{Linear Recursion and Iteration}
\label{Section 1.2.1}

We begin by considering the factorial function, defined by

$$n! = n \cdot (n - 1) \cdot (n - 2) \cdots 3 \cdot 2 \cdot 1.$$

There are many ways to compute factorials.  One way is to make use of the
observation that \( n! \) is equal to \( n \) times \( (n - 1)! \) for any positive
integer \( n \):

$$n! = n \cdot [(n - 1) \cdot (n - 2) \cdots 3 \cdot 2 \cdot 1] = n \cdot (n - 1)!.$$

Thus, we can compute \( n! \) by computing \( (n - 1)! \) and multiplying the
result by \( n \).  If we add the stipulation that 1! is equal to 1, this
observation translates directly into a procedure:

\begin{codeBlock}
>>> def factorial(n):
...     if n == 1:
...         return 1
...     else:
...         return n * factorial(n - 1)
\end{codeBlock}

\noindent
We can use the substitution model of \link{Section 1.1.5} to watch this
procedure in action computing 6!, as shown in \link{Figure 1.3}.

% figure 1.3
\begin{imageFigure}{A linear recursive process for computing 6!.}
  % A linear recursive process for computing 6!, used in Section 1.2.1.
%
% The substitution trace of the recursive factorial, with the arrow that is the
% point of the figure: the process expands as the chain of deferred
% multiplications is built, then contracts as they are performed.
%
% The parentheses are not decoration. Multiplication in Python associates to
% the left, so 6 * 5 * 4 would group as ((6 * 5) * 4) -- a different tree from
% the one the recursion builds, which is 6 * (5 * (4 * ...)). The substitution
% produces the nesting shown.

\begin{tikzpicture}[
  x = 1mm, y = 4.6mm,
  line/.style = {anchor=west, font=\ttfamily\small, inner sep=0pt},
]

\node[line] (l1)  at (0,  0) {factorial(6)};
\node[line]       at (0, -1) {6 * factorial(5)};
\node[line]       at (0, -2) {6 * (5 * factorial(4))};
\node[line]       at (0, -3) {6 * (5 * (4 * factorial(3)))};
\node[line]       at (0, -4) {6 * (5 * (4 * (3 * factorial(2))))};
\node[line] (mid) at (0, -5) {6 * (5 * (4 * (3 * (2 * factorial(1)))))};
\node[line]       at (0, -6) {6 * (5 * (4 * (3 * (2 * 1))))};
\node[line]       at (0, -7) {6 * (5 * (4 * (3 * 2)))};
\node[line]       at (0, -8) {6 * (5 * (4 * 6))};
\node[line]       at (0, -9) {6 * (5 * 24)};
\node[line]       at (0,-10) {6 * 120};
\node[line] (l12) at (0,-11) {720};

% Out to the right as the process expands, back to the left as it contracts.
% The path is placed relative to the widest node rather than by guessed
% coordinates, so it clears the text whatever the font metrics turn out to be.
\draw[-{Triangle[open, length=1.8mm]}, thin]
  ([xshift=2mm] l1.east) -- ([xshift=8mm] l1.east);
\draw[thin]
  ([xshift=8mm] l1.east)
    to[out=0, in=90] ([xshift=14mm] mid.east)
    to[out=270, in=0] ([xshift=8mm] l12.east);
\draw[-{Triangle[open, length=1.8mm]}, thin]
  ([xshift=8mm] l12.east) -- ([xshift=2mm] l12.east);

\end{tikzpicture}

\end{imageFigure}

Now let's take a different perspective on computing factorials.  We could
describe a rule for computing \( n! \) by specifying that we first multiply 1 by
2, then multiply the result by 3, then by 4, and so on until we reach \( n \).
More formally, we maintain a running product, together with a counter that
counts from 1 up to \( n \).  We can describe the computation by saying that the
counter and the product simultaneously change from one step to the next
according to the rule

\begin{syntaxForm}
product ← counter * product
counter ← counter + 1
\end{syntaxForm}

\noindent
and stipulating that \( n! \) is the value of the product when the counter
exceeds \( n \).

Once again, we can recast our description as a procedure for computing
factorials:%
\footnote{%
In a real program we would probably use the block
structure introduced in the last section to hide the definition of
\code{fact-iter}:

\begin{compactCodeBlock}
>>> def factorial(n):
...     def iter(product, counter):
...         if counter > n:
...             return product
...         else:
...             return iter(counter * product, counter + 1)
...     return iter(1, 1)
\end{compactCodeBlock}

We avoided doing this here so as to minimize the number of things to think
about at once.%
}

\begin{codeBlock}
>>> def factorial(n):
...     return fact_iter(1, 1, n)
>>> def fact_iter(product, counter, max_count):
...     if counter > max_count:
...         return product
...     else:
...         return fact_iter(counter * product,
...                          counter + 1,
...                          max_count)
\end{codeBlock}

\noindent
As before, we can use the substitution model to visualize the process of
computing 6!, as shown in \link{Figure 1.4}.

% figure 1.4
\begin{imageFigure}{A linear iterative process for computing 6!.}
  % A linear iterative process for computing 6!, used in Section 1.2.1.
%
% The companion to factorial-recursive-process.tex, and the contrast between
% the two is the argument of the section. There the arrow bulges, because the
% process expands as deferred multiplications accumulate and contracts as they
% are performed. Here it is square: the process neither grows nor shrinks, and
% each line is the same size as the last -- three numbers.
%
% The arrow is drawn with right angles deliberately, for that reason.

\begin{tikzpicture}[
  x = 1mm, y = 4.6mm,
  line/.style = {anchor=west, font=\ttfamily\small, inner sep=0pt},
]

\node[line] (l1) at (0,  0) {factorial(6)};
\node[line] (l2) at (0, -1) {fact\_iter(1, 1, 6)};
\node[line]      at (0, -2) {fact\_iter(1, 2, 6)};
\node[line]      at (0, -3) {fact\_iter(2, 3, 6)};
\node[line]      at (0, -4) {fact\_iter(6, 4, 6)};
\node[line]      at (0, -5) {fact\_iter(24, 5, 6)};
\node[line]      at (0, -6) {fact\_iter(120, 6, 6)};
\node[line] (lw) at (0, -7) {fact\_iter(720, 7, 6)};
\node[line] (l9) at (0, -8) {720};

\draw[-{Triangle[open, length=1.8mm]}, thin]
  ([xshift=2mm] l1.east) -- ([xshift=8mm] l1.east);
\draw[thin]
  ([xshift=8mm] l1.east)
    -- ([xshift=10mm] lw.east |- l1.east)
    -- ([xshift=10mm] lw.east |- l9.east)
    -- ([xshift=8mm] l9.east);
\draw[-{Triangle[open, length=1.8mm]}, thin]
  ([xshift=8mm] l9.east) -- ([xshift=2mm] l9.east);

\end{tikzpicture}

\end{imageFigure}

Compare the two processes.  From one point of view, they seem hardly different
at all.  Both compute the same mathematical function on the same domain, and
each requires a number of steps proportional to \( n \) to compute \( n! \).
Indeed, both processes even carry out the same sequence of multiplications,
obtaining the same sequence of partial products.  On the other hand, when we
consider the ``shapes'' of the two processes, we find that they evolve quite
differently.

Consider the first process. The substitution model reveals a shape of
expansion followed by contraction, indicated by the arrow in \mbox{\link{Figure 1.3}}.
The expansion occurs as the process builds up a chain of \newterm{deferred
operations} (in this case, a chain of multiplications).  The contraction occurs
as the operations are actually performed.  This type of process, characterized
by a chain of deferred operations, is called a \newterm{recursive process}.
Carrying out this process requires that the interpreter keep track of the
operations to be performed later on.  In the computation of \( n! \), the length
of the chain of deferred multiplications, and hence the amount of information
needed to keep track of it, grows linearly with \( n \) (is proportional to
\( n \)), just like the number of steps.  Such a process is called a
\newterm{linear recursive process}.

By contrast, the second process does not grow and shrink.  At each step, all we
need to keep track of, for any \( n \), are the current values of the variables
\code{product}, \code{counter}, and \code{max\_count}.  We call this an
\newterm{iterative process}.  In general, an iterative process is one whose
state can be summarized by a fixed number of \newterm{state variables},
together with a fixed rule that describes how the state variables should be
updated as the process moves from state to state and an (optional) end test
that specifies conditions under which the process should terminate.  In
computing \( n! \), the number of steps required grows linearly with \( n \).  Such
a process is called a \newterm{linear iterative process}.

The contrast between the two processes can be seen in another way.  In the
iterative case, the program variables provide a complete description of the
state of the process at any point.  If we stopped the computation between
steps, all we would need to do to resume the computation is to supply the
interpreter with the values of the three program variables.  Not so with the
recursive process.  In this case there is some additional ``hidden''
information, maintained by the interpreter and not contained in the program
variables, which indicates ``where the process is'' in negotiating the chain of
deferred operations.  The longer the chain, the more information must be
maintained.\footnote{When we discuss the implementation of procedures on
register machines in \link{Chapter 5}, we will see that any iterative process
can be realized ``in hardware'' as a machine that has a fixed set of registers
and no auxiliary memory.  In contrast, realizing a recursive process requires a
machine that uses an auxiliary data structure known as a \newterm{stack}.}

In contrasting iteration and recursion, we must be careful not to confuse the
notion of a recursive \newterm{process} with the notion of a recursive
\newterm{procedure}.  When we describe a procedure as recursive, we are
referring to the syntactic fact that the procedure definition refers (either
directly or indirectly) to the procedure itself.  But when we describe a
process as following a pattern that is, say, linearly recursive, we are
speaking about how the process evolves, not about the syntax of how a procedure
is written.  It may seem disturbing that we refer to a recursive procedure such
as \code{fact\_iter} as generating an iterative process.  However, the process
really is iterative: Its state is captured completely by its three state
variables, and an interpreter need keep track of only three variables in order
to execute the process.

One reason that the distinction between process and procedure may be confusing
is that most implementations of common languages are designed in such a way
that the interpretation of any recursive procedure consumes an amount of memory
that grows with the number of procedure calls, even when the process described
is, in principle, iterative.  As a consequence, these languages can describe
iterative processes only by resorting to special-purpose ``looping constructs''
such as \code{for} and \code{while}.  An implementation that does not share
this defect---that executes an iterative process in constant space even when
the process is described by a recursive procedure---is called
\newterm{tail-recursive}.  With a tail-recursive implementation, iteration can
be expressed using the ordinary procedure call mechanism, so that special
iteration constructs are useful only as syntactic sugar.\footnote{Tail
recursion has long been known as a compiler optimization trick.  A coherent
semantic basis for tail recursion was provided by Carl \link{Hewitt (1977)},
who explained it in terms of the ``message-passing'' model of computation that
we shall discuss in \link{Chapter 3}.  Inspired by this, Gerald Jay Sussman and
Guy Lewis Steele Jr. (see \link{Steele and Sussman 1975}) constructed a
tail-recursive interpreter for Scheme.  Steele later showed how tail recursion
is a consequence of the natural way to compile procedure calls
(\link{Steele 1977}).}

\begin{departure}
Scheme is tail-recursive; the standard requires it.  Python is not, and this is
the first place in the book where the difference costs us something we cannot
argue away.

Run \code{factorial(10000)} with either definition above and Python answers
with a \code{RecursionError}.  For the first that is no surprise---the process
really is recursive, and the chain of deferred multiplications has to be kept
somewhere.  But the second fails in exactly the same way, and it describes an
iterative process.  Its state is three numbers.  Python keeps a frame for the
call regardless, because it does not notice that the call is the last thing
\code{fact\_iter} does, and so a procedure that needs to remember nothing
accumulates ten thousand frames remembering it.

The decision is deliberate rather than an oversight.  Eliminating a tail call
means discarding the frame that made it, and that frame is what a traceback is
built from; Python's authors have preferred debuggable stack traces to
unbounded iteration by procedure call.  It is a real trade and not an obviously
wrong one.  But it means that for us the distinction between a recursive
procedure and an iterative process, which the original can leave to the
interpreter, is one we must do something about.
\end{departure}

\subsubsection*{Deferring the call}

There are two things we can do about it, and they are worth seeing in that
order, because the first keeps the shape of the procedure and the second
explains why it works.

The first is to say explicitly what Python failed to notice.  Instead of
calling \code{fact\_iter} in tail position, the procedure returns a description
of the call to be made next, and something else makes it:

\begin{codeBlock}
>>> from pysicp import TailCall, tail_recursive

>>> @tail_recursive
... def fact_iter(product, counter, max_count):
...     if counter > max_count:
...         return product
...     else:
...         return TailCall(fact_iter,
...                         counter * product,
...                         counter + 1,
...                         max_count)
\end{codeBlock}

\noindent
Two things have changed and nothing else has.  The procedure is marked with
\code{@tail\_recursive}, which is Python's notation for handing a procedure to
another procedure that will run it.\footnote{The line \code{@tail\_recursive}
above a definition is called a \newterm{decorator}.  It means what
\code{fact\_iter = tail\_recursive(fact\_iter)} would mean: the procedure just
defined is passed to \code{tail\_recursive}, and whatever comes back is bound to
the name.  We will make our own in \link{Section 1.3}, once procedures that take
and return procedures are familiar; for now it is enough that it does
something to the procedure it is written above.} And the recursive call is
written \code{TailCall(fact\_iter, ...)} rather than
\code{fact\_iter(...)}---the same procedure and the same arguments, but
described rather than performed.

That is enough.  \code{factorial(10000)} now answers with a number of 35660
digits, and the process runs in constant space, because no call is ever nested
inside another.  The shape of the original is preserved: the procedure is still
recursive, the state is still three variables, and the reader can still see the
iterative process in the text of the definition.

Only a call whose value is returned unchanged may be described this way.  Our
first \code{factorial} multiplies by \code{n} after the call returns, so there
is nothing to defer; that process is recursive and the space it consumes is not
an artefact of the implementation but a property of the process itself.

\subsubsection*{The loop that was there all along}

What does \code{tail\_recursive} do with the description?  It makes the call,
looks at what comes back, and if that is another description, makes that call
too---going round until something that is not a description comes back:

\begin{syntaxForm}
result ← the procedure's value
while result is a TailCall:
    result ← make the call it describes
\end{syntaxForm}

\noindent
Which is a loop.  The iterative process was there in the procedure all along,
and \code{tail\_recursive} does no more than run it as one---exactly the
substitution of \link{Section 1.1.5}, performed by a \code{while} rather than
by nesting calls.  This is why the space is constant: the loop has one
\code{result} and updates it, in the same way the process has three state
variables and updates them.

Seeing that, we can write the loop ourselves and dispense with the machinery:

\begin{codeBlock}
>>> def factorial(n):
...     product, counter = 1, 1
...     while counter <= n:
...         product, counter = counter * product, counter + 1
...     return product
\end{codeBlock}

\noindent
This is what a Python programmer would ordinarily write, and it is the version
to prefer.  The state variables are the same three, updated by the same rule,
tested by the same end test; what has gone is the pretence that a procedure
call was being made at each step.  In a tail-recursive language the two forms
are the same program written twice, and the choice between them is a matter of
taste.  In Python the second is what the first would have been, and writing it
is a transformation the programmer performs rather than one the interpreter
performs for them.

We will use both.  Where the shape of the recursion is the point---and in this
chapter it usually is---we will keep the recursive procedure, with
\code{TailCall} where it is needed to make the space behaviour honest.  Where
the loop is the point, we will write the loop.  And in \link{Chapter 5}, when
we build a machine to run this language, we will implement tail calls in it
properly, and the loop we have just written by hand will be the thing the
machine does mechanically.

\begin{exerciseGroup}
\phantomsection\label{Exercise 1.9}
\exerciseHeading{Exercise 1.9:} Each of the following two
procedures defines a method for adding two positive integers in terms of the
procedures \code{inc}, which increments its argument by 1, and \code{dec},
which decrements its argument by 1.  The original defines these as \code{+},
redefining the operator itself; Python's operators are not names and cannot be
redefined, as we saw in \link{Section 1.1.1}, so they are written here as an
ordinary procedure \code{add}.

\begin{codeBlock}
>>> def add(a, b):
...     return b if a == 0 else inc(add(dec(a), b))

>>> def add(a, b):
...     return b if a == 0 else add(dec(a), inc(b))
\end{codeBlock}

Using the substitution model, illustrate the process generated by each
procedure in evaluating \code{add(4, 5)}.  Are these processes iterative or
recursive?

\phantomsection\label{Exercise 1.10}
\exerciseHeading{Exercise 1.10:} The following procedure computes
a mathematical function called Ackermann's function.

\begin{codeBlock}
>>> def A(x, y):
...     if y == 0:
...         return 0
...     elif x == 0:
...         return 2 * y
...     elif y == 1:
...         return 2
...     else:
...         return A(x - 1, A(x, y - 1))
\end{codeBlock}

What are the values of the following expressions?

\begin{codeBlock}
A(1, 10)
A(2, 4)
A(3, 3)
\end{codeBlock}

Consider the following procedures, where \code{A} is the procedure
defined above:

\begin{codeBlock}
>>> def f(n):
...     return A(0, n)
>>> def g(n):
...     return A(1, n)
>>> def h(n):
...     return A(2, n)
>>> def k(n):
...     return 5 * n * n
\end{codeBlock}

Give concise mathematical definitions for the functions computed by the
procedures \code{f}, \code{g}, and \code{h} for positive integer values of
\( n \).  For example, \code{k(n)} computes \( 5n^2 \).
\end{exerciseGroup}

